\section{The presentation on scan 336}

We work on the big fppf site of the indicated base. A constant group
means the associated locally constant sheaf, not a diagonalizable group.
In this section alone, $s$ denotes the earlier lift $s_-$ of the
conventions table, because it is the lift used on scan 336.
We write $x(t)=\left(\begin{smallmatrix}1&t\\0&1\end{smallmatrix}\right)$,
$s=\left(\begin{smallmatrix}0&-1\\1&0\end{smallmatrix}\right)$, and
$z=-I$. In the normalization of \S48.VI, scan 336, displayed folio 335,
the matrices are
\[
 \sigma=s,\qquad \rho=\begin{pmatrix}0&1\\-1&1\end{pmatrix},
 \qquad \varepsilon _0(t)=x(-t).
\]
Thus $\rho^{-1}\sigma\varepsilon _0(1)=\rho^3=\sigma^2=z$.
The manuscript really says $\mathbb Z/2$, $\mathbb Z/3$, $\mathbb Z/4$,
and $\mathbb Z/6$: the photograph was inspected. Replacing these by
$\mu_n$ changes the question.

\begin{statement}{C.1}{Proved}{The literal sheaf presentation}
Put $\Gamma=\mathrm{SL}_2(\mathbb Z)$, $\overline\Gamma=
\mathrm{PSL}_2(\mathbb Z)$ and let $\mathbb Z\to\Gamma$ send $1$ to
$x(1)$. The two sheaves specified by scan 336 are
\[
 A_{\mathrm{SL}}=\underline\Gamma*_{\underline{\mathbb Z}}\mathbb G_a,
 \qquad A_{\mathrm P}=
 \underline{\overline\Gamma}*_{\underline{\mathbb Z}}\mathbb G_a.       \tag{C.1}
\]
Here the asterisk is a pushout in sheaves of groups, and the other map
from $\underline{\mathbb Z}$ sends $n$ to $n\cdot1$.
They have canonical comparison homomorphisms
$A_{\mathrm{SL}}\to\mathrm{SL}_2$ and $A_{\mathrm P}\to\mathrm{PGL}_2$.
The same generators and relation that the manuscript assigns to
$\mathrm{GL}_2$ give $A_{\mathrm{SL}}\to\mathrm{SL}_2\hookrightarrow
\mathrm{GL}_2$.
\end{statement}
\begin{proof}
The first two finite factors in the special-linear version give
$C_4*_{C_2}C_6$, with the common element
$\sigma^2=\rho^3$ of order two. The additional relation (63) identifies
$\varepsilon _0(1)$ with $\sigma^{-1}\rho z$. Eliminating this generator
gives precisely the pushout (C.1); changing $t$ to $-t$ gives our $x$.
In the projective version one starts with $C_2*C_3$ and eliminates
$\varepsilon _0(1)$ using (62). The identification of these discrete
groups with $\Gamma$ and $\overline\Gamma$ is proved next.
All displayed matrices have determinant one. Consequently the asserted
$\mathrm{GL}_2$ presentation cannot include matrices with arbitrary
determinant. This conclusion uses the actual matrices, independently
of any doubtful word surrounding the formulas.
\end{proof}

\begin{statement}{C.2}{Known}{The discrete integral amalgam and its reduction}
The comparison maps
\[
 C_2*C_3\xrightarrow{\sim}\mathrm{PSL}_2(\mathbb Z),\qquad
 C_4*_{C_2}C_6\xrightarrow{\sim}\mathrm{SL}_2(\mathbb Z)
\]
are isomorphisms of discrete groups. This is the conclusion of Conrad,
Theorem C.1 and Theorem 1.1, in the present normalization
\cite{ConradSL2Z}; the following proof gives the reduction and the
central lift explicitly.
\end{statement}
\begin{proof}
For $\left(\begin{smallmatrix}a&b\\c&d\end{smallmatrix}\right)$ of
determinant one over $\mathbb Z$, left multiplication by $x(n)$ adds
$n$ times the second row to the first and left multiplication by $s$
interchanges the rows with a sign. The Euclidean algorithm reduces
$(a,c)$ to $(\pm1,0)$. The remaining matrix is a power of $x(1)$,
possibly multiplied by $s^2=-I$. Hence $s,x(1)$ generate.

For injectivity in the projective group, cyclically reduce an alleged
nontrivial kernel word in $C_2*C_3$. A word in a single factor has
nonidentity image. Otherwise cyclically rotate it to the form
$(s\rho^{i_1})\cdots(s\rho^{i_m})$ with $i_j=1$ or $2$.
Direct multiplication gives $s\rho=x(-1)$ and $s\rho^2=y(-1)$.
Conjugation by $\operatorname{diag}(-1,1)$ makes these blocks
$x(1),y(1)$. Starting with $I$, multiplication by either positive
block strictly increases the sum of the nonnegative entries, because
both rows and columns of an invertible such matrix have positive sum.
A nonempty product therefore has entry sum greater than two and
cannot be $I$; nonnegative entries also exclude $-I$.
This is Conrad's positive-matrix argument, with the change of
normalization exhibited. It proves injectivity without a geometric
fundamental-domain assumption.

In $C_4*_{C_2}C_6$ the common involution is central, because it is
central in each factor. Quotienting by it gives $C_2*C_3$. The preceding
injectivity shows that the kernel of the special-linear comparison is
contained in this order-two subgroup; its nonidentity element maps
to $-I\ne I$. The kernel is therefore trivial. This also proves that
the parabolic subgroup used in (C.1) is an embedded copy of $\mathbb Z$.
\end{proof}

\begin{statement}{C.3}{Proved}{Images and cokernels in the fppf topology}
Over every base, $A_{\mathrm{SL}}\to\mathrm{SL}_2$ is an epimorphism
even for the Zariski topology, and $A_{\mathrm P}\to\mathrm{PGL}_2$
is an fppf epimorphism. The image of the literal general-linear map is
$\mathrm{SL}_2$ and its sheaf cokernel is $\mathbb G_m$, by determinant.
For a local ring $R$, the subgroup of $\mathrm{SL}_2(R)$ generated by
$s$ and the $x(t)$ is all of $\mathrm{SL}_2(R)$. For a local ring, or
for $R=\mathbb Z,\mathbb Z[1/p]$, the subgroup they generate in
$\mathrm{PGL}_2(R)$ is the kernel of
\[
 \det\bmod (R^\times)^2:\mathrm{PGL}_2(R)\longrightarrow
 R^\times/(R^\times)^2.                                             \tag{C.2}
\]
This last assertion concerns the subgroup generated by the displayed
global points, not an interchange of evaluation with a sheaf pushout.
\end{statement}
\begin{proof}
Put $y(v)=s x(-v)s^{-1}$ and
\[
 w(u)=x(u)sx(u^{-1})s^{-1}x(u),\qquad h(u)=w(u)s
       =\operatorname{diag}(u,u^{-1})\quad(u\in R^\times).
\]
These identities hold over $\mathbb Z[u,u^{-1}]$. If $a$ is a unit,
then a determinant-one matrix factors as
\[
 \begin{pmatrix}a&b\\c&d\end{pmatrix}
 =y(c/a)h(a)x(b/a).                                                 \tag{C.3}
\]
If $c$ is a unit, apply the same formula to $s^{-1}g$.
The two open sets $D(a)$ and $D(c)$ cover the base since $ad-bc=1$.
This proves the Zariski epimorphism and generation over local rings.
Over $\mathbb Z$ use C.2; over $\mathbb Z[1/p]$ clear denominators
and perform the integral Euclidean algorithm. The resulting gcd is
a power of $p$ and hence a unit, so (C.3) finishes the reduction.

A projective-linear point is fppf locally represented by a matrix in
$\mathrm{GL}_2$. Adjoining a square root of its unit determinant is
a finite free faithfully flat extension, including in characteristic
two. Scalar multiplication then makes the determinant one. This proves
the projective fppf epimorphism. If $\operatorname{Pic}(R)=0$, a
projective-linear point has a matrix representative; it has a
determinant-one representative exactly when that determinant is a
square. All rings specified in the statement have trivial Picard group.
Finally determinant gives the split exact sequence
$1\to\mathrm{SL}_2\to\mathrm{GL}_2\to\mathbb G_m\to1$,
with section $d\mapsto\operatorname{diag}(d,1)$.
\end{proof}

\begin{statement}{C.4}{Refuted}{The proposed presentation fails over every nonempty base}
Neither comparison in C.1 is an isomorphism over any nonempty base
scheme. Their kernels are noncentral after a field extension of every
geometric fibre. In particular there is no nonempty open subset of
$\operatorname{Spec}\mathbb Z$, including $\operatorname{Spec}
\mathbb Z[1/6]$, over which the literal presentation is correct.
The kernel is therefore not the central $K_2$ kernel of a Steinberg
presentation. Lead 1, in the form proposed in the prompt, is false.
\end{statement}
\begin{proof}
We give a word surviving fppf sheafification; this is stronger than
merely finding a relation absent from a formal list. In characteristic
zero put $H=\mathbb Z$ and let $\Gamma_0$ be $\Gamma$ or
$\overline\Gamma$. In characteristic $p$, put
\[
 \Gamma_p=\Gamma_0/\langle\!\langle x(1)^p\rangle\!\rangle,
 \qquad H=C_p.
\]
There is a homomorphism from $\Gamma_p$ to the appropriate matrix
group over $\mathbb F_p$. Its restriction to $H$ is injective, and
$s\notin H$, as seen from the lower-left entry. On a connected
algebra in that characteristic, the presheaf underlying (C.1) is the
ordinary amalgam $\Gamma_p*_{H}(R,+)$; in characteristic zero it is
$\Gamma_0*_{\mathbb Z}(R,+)$. This follows by first imposing exactly
the kernel of $\mathbb Z\to(R,+)$ and then taking the ordinary pushout.

Let $k$ be any field and work in $k(T,U,V)$, with independent variables.
The matrix word
\[
 K=h(T)x(U)h(T)^{-1}x(-T^2U)
\]
is the identity, since $h(T)=\operatorname{diag}(T,T^{-1})$.
Conjugating it by $x(-T)$ gives the following alternating word $L$:
\[
 \begin{split}
 L={}&s x(T^{-1})s^{-1}x(T)s x(U)s^{-1}x(-T)\\
     &\hspace{15mm}{}s x(-T^{-1})s^{-1}x(-T^2U).
 \end{split}
 \tag{C.4}
\]
Every $s^{\pm1}$ lies outside $H$, and every displayed additive
parameter lies outside $H$. Thus (C.4) is a cyclically reduced word
of length twelve in the amalgam, and is nonidentity. The commutator
$[L,x(V)]$ has length twenty-four after cancellation of the adjacent
additive parameters: its central additive parameter is $V$, and all
other additive syllables are still nonconstant. Hence it too is
nonidentity, proving noncentrality.

For completeness, the normal-form fact used here is elementary.
Choose representatives for right cosets of the common subgroup in
each factor. Move common-subgroup elements to the right successively;
induction on the number of alternating factors gives a unique reduced
representative. A reduced alternating word of positive length cannot
represent the identity. This proof also proves injectivity of the two
factors. It applies to our two factors because the maps from $H$ are
injective, as checked above.

No faithfully flat extension can invalidate this argument. A nonzero
algebra over $k(T,U,V)$ contains that field injectively. Its displayed
parameters therefore remain outside the constant subgroup $H$ on
every nonempty connected piece. The same reduced word remains
nonidentity on every fppf cover. Thus it is nonidentity in the sheaf
pushout, and so is its commutator. Every nonempty base has a field
point; base change preserves the sheaf presentation. This proves the
assertion for every base.

There are smaller rational witnesses. In the special-linear case
$[x(1/2),z]$ is a reduced nonidentity word but the identity matrix.
In the projective case $w(2/3)^2$ is a nonidentity reduced amalgam
word but has matrix $-I$. The symbolic program verifies both matrix
identities. Its integral calculation of (C.4) verifies the same
identity simultaneously in characteristics $2,3,5$.
\end{proof}

The distinction from the $K$-theoretic literature is substantive.
Brylinski--Deligne define central extensions on the big Zariski site
and classify those extensions, not the much larger source (C.1)
\cite[Introduction and Theorem 7.2]{BrylinskiDeligne2001C}.
Moreover their introduction explicitly distinguishes type $C_n$,
including $\mathrm{SL}_2$, in the description of universal central
extensions. Our noncentral kernel already prevents application of
that classification to the literal manuscript presentation.

\begin{statement}{C.5}{Proved}{An explicit normal generating family for the entire kernel}
The kernel of $A_{\mathrm{SL}}\to\mathrm{SL}_2$ is the normal subsheaf
generated by the following four families of words. All symbols below
denote words in $s$ and $x(t)$, so this is an explicit presentation.
For $u$ invertible put
\[
 N_0(v,u,t)=y(v)h(u)x(t),\qquad N_1(v,u,t)=sN_0(v,u,t),
 \qquad D=1+u^2vt.
\]
For $i=0,1$ impose
\begin{align}
 N_i(v,u,t)s&=N_i\bigl(v+u^{-2}t^{-1},ut,-t^{-1}\bigr)
                   &&(t\text{ invertible}),                    \tag{C.5}\\
 N_i(v,u,t)s&=N_{1-i}\bigl(-u^2t/D,(-1)^iD/u,-u^2v/D\bigr)
                   &&(D\text{ invertible}).                    \tag{C.6}
\end{align}
These relations are valid over $\mathbb Z$, with no primes inverted.
They identify the entire kernel as a normal subsheaf, rather than only
exhibiting some of its elements.
\end{statement}
\begin{proof}
Multiplication gives
\[
 N_0(v,u,t)=\begin{pmatrix}u&ut\\uv&uvt+u^{-1}\end{pmatrix},\qquad
 N_1(v,u,t)=\begin{pmatrix}-uv&-uvt-u^{-1}\\u&ut\end{pmatrix}.
\]
Substitution proves (C.5)--(C.6). The attached program independently
reduces each of the four matrix differences to zero over the
corresponding universal Laurent ring. The open sets $D(t)$ and $D(D)$
cover, because $D-u^2vt=1$.

Let $Q$ be the sheaf quotient by these relators. A local word in
$s,x(t)$ reduces, after a Zariski cover, to one of the forms $N_i$.
Indeed the identity is $N_0(0,1,0)$: $h(1)=1$ is already an identity
in the discrete factor $\Gamma$. Right multiplication by $x(a)$
changes $t$ to $t+a$. Right multiplication by $s$ is handled by
(C.5)--(C.6) on their covering open sets. Inverses cause no problem:
$x(a)^{-1}=x(-a)$ and $s^{-1}=s^3$. Induction on word length proves
the asserted local reduction.

If a section of $Q$ maps to the identity matrix, reduce it locally.
It cannot have form $N_1$, whose lower-left entry is a unit.
In form $N_0$ its entries force $u=1$, $v=t=0$, so it is the identity
in $Q$. Thus $Q\to\mathrm{SL}_2$ is a monomorphism of sheaves;
surjectivity is C.3. The quotient is exactly $\mathrm{SL}_2$ and
the stated normal subsheaf is exactly the kernel. Notice that the
proof is valid over rings with nilpotents as well as fields.
\end{proof}

\begin{statement}{C.6}{Proved}{The corresponding presentations for projective and general linear groups}
For $A_{\mathrm P}$, add the images of (C.5)--(C.6) and also
\[
 h(\zeta)=1\quad\text{for every }\zeta\in\mu_2.                  \tag{C.7}
\]
Their normal closure is exactly $\ker(A_{\mathrm P}\to\mathrm{PGL}_2)$.
Over $\mathbb Z[1/2]$, (C.7) follows locally from $s^2=1$ after
(C.5)--(C.6); over characteristic two it cannot be omitted.

A presentation of $\mathrm{GL}_2$ is obtained from the corrected
special-linear presentation by adjoining a copy $m(\mathbb G_m)$
and imposing
\[
 m(d)x(t)m(d)^{-1}=x(dt),\qquad
 m(d)s m(d)^{-1}=h(d)s.                                           \tag{C.8}
\]
Here $m(d)$ maps to $\operatorname{diag}(d,1)$.
\end{statement}
\begin{proof}
The source $A_{\mathrm P}$ is the quotient of $A_{\mathrm{SL}}$ by
the normal subsheaf generated by the constant element $z=s^2$.
After the Gaussian relations, quotienting by all scalar sections
$\zeta I$, $\zeta^2=1$, gives the fppf quotient
$\mathrm{SL}_2/\mu_2=\mathrm{PGL}_2$, by the same square-root
argument as C.3. Since $h(\zeta)=\zeta I$ in $\mathrm{SL}_2$,
(C.7) imposes exactly this quotient.
When two is invertible, the roots of $X^2-1$ are fppf locally
$1,-1$, so the constant relation suffices locally. In
$R=\mathbb F_2[e]/(e^2)$, however, $\zeta=1+e$ satisfies
$\zeta^2=1$ and $\zeta\ne1$. The scalar $h(\zeta)=\zeta I$
is nonidentity in $\mathrm{SL}_2(R)$ and identity in
$\mathrm{PGL}_2(R)$. This finite counterexample is checked in the
program. In characteristic two the constant $-I$ is already $I$;
the extra scalar cannot be removed by that constant relation.

For (C.8), conjugation by $\operatorname{diag}(d,1)$ has exactly
the two stated effects. Since $x(t)$ and $s$ generate as a sheaf,
these relations impose its complete conjugation action on
$\mathrm{SL}_2$. The presented sheaf is therefore the semidirect
product $\mathrm{SL}_2\rtimes\mathbb G_m$ for that action.
Every invertible matrix $g$ has the unique expression
$(g\operatorname{diag}(\det(g)^{-1},1))\operatorname{diag}(\det(g),1)$,
which proves its isomorphism with $\mathrm{GL}_2$.
\end{proof}

\begin{statement}{C.7}{Proved}{Executed point presentations and their kernels}
Let $P_{\mathrm{SL}}(R)$ and $P_{\mathrm P}(R)$ mean the ordinary
abstract presentations obtained by evaluating the generators and
relations before sheafification. These are distinguished from
$A_{\mathrm{SL}}(R),A_{\mathrm P}(R)$. For $R=\mathbb Z/n$,
\[
 \begin{split}
 P_{\mathrm P}(R)&=\langle s,t\mid s^2=(st)^3=t^n=1\rangle,\\
 P_{\mathrm{SL}}(R)&=\langle s,t\mid s^4=1,
                         (st)^3=s^2,\ t^n=1\rangle.
 \end{split}
 \tag{C.9}
\]
For $n=2,3,4,5$ their respective orders are
$(6,12,24,60)$ and $(6,24,48,120)$. In each case the comparison
is an isomorphism onto the subgroup generated by the corresponding
matrices. For $R=\mathbb F_4$, both point presentations are
$S_3*_{C_2}(C_2\times C_2)$, and the kernel of the comparison
to $\mathrm{SL}_2(\mathbb F_4)$ is a free group of rank six.
For $R=\mathbb F_7$ and $\mathbb Z/9$, the kernels of both
point comparisons onto their generated matrix images are surface
groups of genera three and ten, respectively.
\end{statement}
\begin{proof}
For a cyclic additive ring, every $x(a)$ is a power of $x(1)$.
Eliminating those generators gives (C.9), using the equivalent
standard discrete generators $s,x(1)$ from C.2. The central element
$s^2$ in the special-linear presentation has order at most two;
its quotient is the projective presentation.

The program performs Todd--Coxeter enumeration for the trivial
subgroup for $n=2,3,4,5$ in both presentations. It saves every
coset table, checks that each generator column is a permutation
with the indicated inverse, and checks every relator at every
coset. The eight completed enumerations have the orders stated.
For an independent geometric description of the projective upper
bounds, they are the rotational triangle groups of the spherical
triangles $(2,3,n)$, with orders $6,12,24,60$; the enumeration does
not rely on this description. The complete matrix enumerations
give precisely the same orders for the images. A surjective map
between these finite groups of the same order is injective.

In characteristic two the discrete quotient $\Gamma_2$ has order
six, as the enumeration for $n=2$ proves, and is $S_3$; the prime
additive subgroup has order two. The pushout with
$(\mathbb F_4,+)=C_2^2$ is therefore the asserted amalgam.
Its two factors inject into $\mathrm{SL}_2(\mathbb F_4)$, which
has order sixty. The kernel acts freely on the amalgam tree.
The quotient graph has $60/6+60/4=25$ vertices and $60/2=30$
edges; it is connected. Its fundamental group is free of rank
$30-25+1=6$. The tree construction is explicit: vertices are
cosets of the two factors and edges are cosets of $C_2$; normal
forms prove connectedness and absence of circuits. A group acting
freely on a tree identifies with the fundamental group of its
quotient graph by lifting paths. This supplies the required reduction
to the elementary graph calculation.

For $n=7,9$, let $Q$ be the projective matrix image, of order
$168,324$ respectively. The three elements $s,st,t$ have exact
orders $2,3,n$. The regular permutation action of $Q$ constructs
an oriented surface as follows. Take $|Q|$ edges, with cyclic
orders at black and white vertices prescribed by the permutations
of orders two and three. Glue in the faces prescribed by the
permutation of order $n$. This produces a connected orientable
surface because the generators act transitively; its numbers of
vertices, edges and faces are $|Q|/2+|Q|/3$, $|Q|$, and $|Q|/n$.
The exact orders ensure that no vertex or face has residual
branching. The path-lifting construction, or van Kampen applied
to the glued polygons, identifies its fundamental group with
the kernel of the triangle presentation in (C.9): the only
local loops killed upstairs are precisely $s^2,(st)^3,t^n$.
Its Euler characteristic is
\[
 |Q|(1/2+1/3+1/n-1)=-4,-18,
\]
so its genus is three or ten. In particular both presentations
are infinite, a fact not inferred from failure of enumeration.
For the special-linear comparison, $s^2$ maps to $-I\ne I$.
Projection identifies its kernel with the projective kernel:
any lift of a projective-kernel element maps to $I$ or $-I$,
and multiplication by $s^2$ selects the unique lift mapping to $I$.
\end{proof}

The complete matrix counts are as follows. ``Image'' means the
subgroup generated by the global points of the specified factors.
For every row the special-linear image is all of $\mathrm{SL}_2(R)$.
\begin{center}
\begin{tabular}{lrrrrr}
\hline
$R$&$|\mathrm{SL}_2|$&$|\mathrm{GL}_2|$&$|\mathrm{PGL}_2|$
 &$|\text{projective image}|$&projective index\\ \hline
$\mathbb F_2$&6&6&6&6&1\\
$\mathbb F_3$&24&48&24&12&2\\
$\mathbb F_4$&60&180&60&60&1\\
$\mathbb F_5$&120&480&120&60&2\\
$\mathbb F_7$&336&2016&336&168&2\\
$\mathbb Z/4$&48&96&48&24&2\\
$\mathbb Z/9$&648&3888&648&324&2\\ \hline
\end{tabular}
\end{center}
All entries are obtained by enumerating all four matrix coordinates,
testing determinants, and closing the stated generators under
multiplication. Projective equivalence is tested by all unit scalar
multiples. The predicted index in (C.2) is verified independently.

\begin{statement}{C.8}{Proved}{Sheafification changes points over finite rings}
For $q=p^a$, let $\Gamma_p$ denote the appropriate discrete quotient
used in C.4. There is an exact description
\[
 A(\mathbb F_q)=
 \left(\Gamma_p*_{C_p}(\overline{\mathbb F}_p,+)\right)^{{\rm Frob}_q}.
                                                                    \tag{C.10}
\]
Frobenius fixes the discrete factor and acts by $u\mapsto u^q$
on the additive factor. In particular both sheaves have
$A(\mathbb F_2)=S_3$ and
$A(\mathbb F_4)=S_3*_{C_2}(C_2\times C_2)$.
In contrast, $A_{\mathrm{SL}}(\mathbb F_3)$ is infinite although
$P_{\mathrm{SL}}(\mathbb F_3)$ has order twenty-four.
Both $A_{\mathrm{SL}}(\mathbb Z/4)$ and $A_{\mathrm P}(\mathbb Z/4)$
are infinite, although their point presentations have orders
forty-eight and twenty-four. There is also an explicit additional
infinite-order descent kernel element in
$A_{\mathrm{SL}}(\mathbb Z/9)$.
\end{statement}
\begin{proof}
Every nonempty fppf cover of an algebraically closed field has a
rational point: a nonempty scheme of finite presentation over that
field has a closed point with that residue field. Thus sheafification
does not change the value of the presentation presheaf at
$\overline{\mathbb F}_p$. This gives the amalgam on the right of
(C.10). Any word uses finitely many algebraic parameters, so it is
defined over a finite extension. Equality of reduced words is also
detected over a finite extension. Finite Galois descent now identifies
sections over $\mathbb F_q$ with the Frobenius-invariant words.
This proves (C.10), including both injectivity and surjectivity.

For $p=2$, $\Gamma_2=S_3$ and its subgroup $C_2$ is malnormal:
$C_2\cap gC_2g^{-1}=1$ when $g\notin C_2$. The normal-form
comparison of a word and its Frobenius transform allows only
insertions of elements of $C_2$ between adjacent syllables. At
each discrete syllable outside $C_2$, malnormality forces both
adjacent inserted elements to be the identity. Thus every additive
parameter in a reduced invariant word is itself Frobenius-fixed.
The remaining cases, a word in just one factor, give the same
conclusion. This proves the two asserted characteristic-two values.

For the contrasting characteristic-three example, take
$\alpha\in\mathbb F_{27}$ with $\alpha^3-\alpha=1$.
The cubic $X^3-X-1$ has no root in $\mathbb F_3$ and is
irreducible. In the discrete factor $\Gamma_3=\mathrm{SL}_2(\mathbb F_3)$,
$z=-I$ is a nonidentity central involution outside $C_3$.
The word
\[
 b=x(\alpha)z x(-\alpha)
\]
is Frobenius-fixed, since $\alpha^3=\alpha+1$ and $x(1)$
commutes with $z$ in the discrete factor. Hence it descends to
$A_{\mathrm{SL}}(\mathbb F_3)$. The word $bz$ is cyclically
reduced of length four, so $(bz)^m$ has length $4m$ for every
$m>0$. It has infinite order and maps to the identity matrix.
The program verifies $\alpha^3=\alpha+1$ by exact arithmetic
in $\mathbb F_3[X]/(X^3-X-1)$. Thus the finite coset tables
in C.7 must not be described as the values of the sheaf pushout.

For the four-element base ring put
$R=\mathbb Z/4$ and $S=R[b]/(b^2+b+1)$.
This is a finite etale Galois cover with involution
$F(b)=-b-1$. The element $a=2b$ satisfies $F(a)=a+2$,
$2a=0$, and $a\notin R$. In each discrete point factor
from C.7 put $g=y(2)$. Its matrix shows $g\notin\langle x(1)\rangle$
and $gx(2)=x(2)g$. These are identities in the discrete point
factor because C.7 proves its isomorphism with its matrix image.
Consequently $x(a)g x(-a)$ is $F$-invariant and descends.
The word
\[
 x(a)g x(-a)g^{-1}
\]
is cyclically reduced of length four and has infinite order.
Its matrix is the identity: $2a=0$ gives $x(a)y(2)=y(2)x(a)$.
The nontriviality persists on fppf covers. Indeed every nonzero
flat algebra over the local Artinian ring $S$ is faithfully flat
and hence contains $S$ injectively; the parameter $a$ cannot
become an element of the embedded constant subgroup $\mathbb Z/4$
on any nonempty component. The reduced-word argument therefore
applies just as in C.4.

Finally take $R=\mathbb Z/9$ and
$S=R[b]/(b^3-b-1)$. The derivative is a unit and the reduction
of the polynomial is irreducible, so this is a finite etale
Galois cover of degree three. Its Frobenius generator is
$F(b)=1+4b+3b^2$; direct calculation verifies the polynomial
relation and $F^3(b)=b$. Now $a=3b$ satisfies $F(a)=a+3$.
With $z=s^2$ in the special-linear discrete factor,
$x(a)z x(-a)z$ descends, maps to the identity matrix, and is
cyclically reduced of length four. Its powers remain reduced
over every fppf cover for the same local-Artinian argument.
All ring identities in these two examples are executed exactly
in the supplied program; no appeal to a numerical approximation
or to unspecified class field theory is involved.
\end{proof}

\begin{statement}{C.9}{Known}{The two requested localized integral amalgams}
For each prime $p$, in particular $p=2,3$, put
$D_p=\operatorname{diag}(1,p)$ and
$\Gamma_0(p)=\{g\in\mathrm{SL}_2(\mathbb Z):c(g)\equiv0\pmod p\}$.
There is an isomorphism of discrete groups
\[
 \mathrm{SL}_2(\mathbb Z[1/p])=
 \mathrm{SL}_2(\mathbb Z)*_{\Gamma_0(p)}
 D_p\mathrm{SL}_2(\mathbb Z)D_p^{-1}.                              \tag{C.11}
\]
The edge subgroup has index $p+1$ in both vertex groups.
This is the arithmetic application of the tree-amalgam theorem
of Serre \cite[\S1.2]{Serre1969C}; we prove the required arithmetic
reduction here. It is not an assertion that (C.1) becomes an
isomorphism after localizing the base.
\end{statement}
\begin{proof}
Use the tree whose vertices are homothety classes of
$\mathbb Z_p$-lattices in $\mathbb Q_p^2$, with adjacent lattices
of index $p$. The elementary-divisor theorem shows that this
graph is connected and that there is a unique nonbacktracking
path between any two vertices: write their relative elementary
divisors and change the difference one step at a time. It is
therefore a tree. The determinant-one group preserves the parity
of that difference and has no edge inversions.

The stabilizer in $G=\mathrm{SL}_2(\mathbb Z[1/p])$ of the
standard lattice class is $H_0=\mathrm{SL}_2(\mathbb Z)$:
determinant one forces a stabilizing homothety to have valuation
zero, and $\mathbb Z[1/p]\cap\mathbb Z_p=\mathbb Z$.
The adjacent vertex $D_p\mathbb Z_p^2$ has stabilizer
$H_1=D_pH_0D_p^{-1}$. Their intersection is exactly
$\Gamma_0(p)$, by multiplying out the matrices.

The neighbors of the standard vertex are the $p+1$ lines in
$\mathbb F_p^2$. Reduction of $H_0$ is onto
$\mathrm{SL}_2(\mathbb F_p)$, since the latter is generated
by elementary matrices and these lift integrally. This group
acts transitively on those lines. The conjugate statement holds
at the other vertex. By following paths in the tree, the group
generated by $H_0,H_1$ is transitive on edges. Given $g\in G$,
some such product has the same effect on the base edge; their
quotient fixes the edge and lies in $H_0\cap H_1$. Thus these
two subgroups generate $G$ and a single edge is a fundamental
domain. Finally a reduced alternating word in the two stabilizers
moves the base edge along a nonbacktracking path, so cannot be
the identity. This proves the amalgam and the index assertion.
It is the direct proof of the comparison required to apply
Serre's theorem, rather than a citation without reduction.
\end{proof}

\begin{statement}{C.10}{Proved}{The correction over the 2-adic and 3-adic bases}
Over $\operatorname{Spec}\mathbb Z_2$ and
$\operatorname{Spec}\mathbb Z_3$, the corrected special-linear
presentation is exactly C.5, the projective presentation is
exactly C.6, and the general-linear presentation is (C.8) together
with C.5. The literal kernels remain noncentral. No correction
consisting only of replacing a non-etale $\mu_4$ or $\mu_6$
can address the manuscript's question: those are not its factors.
\end{statement}
\begin{proof}
All defining identities and their local-normal-form proofs are
over $\mathbb Z$, and so remain valid under these base changes.
C.4 applies after passage to the generic field, and also gives
independent witnesses in each characteristic. The constant cyclic
schemes in scan 336 are etale in every characteristic; a cyclic
map is not necessarily an embedding. For example $C_4\to
\mathrm{SL}_{2,\mathbb F_2}$ sending a generator to $s$ has
kernel $C_2$, since $s^2=I$ there. The same central collapse
occurs for the $C_6$ map. The actual infinitesimal projective
correction at two is (C.7), exhibited over the four-element
ring in C.6. At three, $\mu_2$ is etale and that extra
infinitesimal phenomenon is absent; the noncentral Gaussian
relations are still necessary.
\end{proof}

\begin{statement}{C.11}{Open}{The remaining evaluation problem in Part C(d)}
The sheaf kernels and corrected presentations have been determined
in C.3--C.6 and C.10. We do not provide a finite or otherwise
simplified group presentation of the Frobenius fixed subgroup
in (C.10) for $A_{\mathrm{SL}}(\mathbb F_3)$, nor a corresponding
complete descent computation of $A_{\mathrm{SL}}(\mathbb Z/4)$,
$A_{\mathrm P}(\mathbb Z/4)$ and the two values over
$\mathbb Z/9$. Thus Part C(d) is partly settled: the raw
point presentations (C.9) have been computed, while the
complete descent calculations for these actual sheaf values remain.
\end{statement}
\begin{proof}[Precise partial result, computation, and obstruction]
The strongest explicit description obtained for the smallest
unresolved finite-field value is
\[
 A_{\mathrm{SL}}(\mathbb F_3)=
 \bigl(\mathrm{SL}_2(\mathbb F_3)*_{C_3}
             (\overline{\mathbb F}_3,+)\bigr)^{u\mapsto u^3}.
\]
C.8 proves this formula and constructs an infinite cyclic
subgroup of its matrix kernel. C.7 computes and completely
certifies the smaller order-twenty-four point presentation;
the Artin--Schreier calculation proves why this computation
does not settle the sheaf value. C.8 does settle both
characteristic-two field values requested in the prompt.

The missing step is an explicit presentation of the fixed
subgroup displayed above, and the analogous descent of
normal forms over the two nonreduced finite rings. Equivalently,
one must classify the compatible changes of reduced amalgam
syllables under finite etale descent, including the nontrivial
$C_3$-torsor and the two unramified ring examples exhibited in C.8;
imposing only relations among
global additive parameters is insufficient. The first missing
case in increasing finite-field size is the indicated
special-linear value over $\mathbb F_3$. No claim of
undecidability or of a previously established open problem
is intended: this is an exact boundary of the present calculation.
\end{proof}

The theorem answering the substantive sheaf question is therefore
C.3--C.6: the generators do generate in the stated topology,
their one displayed relation does not present the group over
any nonempty base, their full kernel is the explicit noncentral
normal relator sheaf, and the Gaussian relations supply a
presentation valid integrally. The remaining point evaluation
question is kept separately in C.11.
